I have a few foundational texts that I would like to lean heavily on

\begin{itemize}
    \item The GDL bible, \href{https://arxiv.org/pdf/2104.13478}
    \item Yuan-Sen's stats "book",  \href{https://arxiv.org/pdf/2506.12230} and ARAA review \href{https://arxiv.org/pdf/2510.10713}
    \item Previous Graph reviews, like this general review \href{https://arxiv.org/pdf/1812.08434}, and this particle physics GNN review \href{https://arxiv.org/pdf/2007.13681}.
\end{itemize}

I have some things that I would like to emphasize

\begin{itemize}
    \item GNNs as a way to avoid Theory noise and being able to transfer better model
    \item The natural intuition that all physics-trained researchers have with Graph Neural Networks. 
    \item Machine-learning in the different N regimes (super high N ($>10^8$) for transformers, low N ($<10^3$) for e.g., random forests, medium N training ($10^4-10^7$), most relevant for statistical astronomy, perfect for us!!
    \item A lot of what we learn in astronomy we learn from the way a given object relates to other objects. Be it transiting exoplanets, galaxy mergers, large-scale structure, gravitational wave events, neutrino detection. Some of these are geometrically simple, while some contain highly non-trivial geometries.
    \item Similarly, the section on when NOT to use a graph neural network should be clear
    \item Weak geometric constraints, a la Andrew Gordon Wilson, as a new/future idea
    \item Domain Adaptation (Natalí's paper)
    
\end{itemize}

Things to remember

\begin{itemize}
    \item Do a table of symmetries relevant to different astrophysics subfields (check if Yuan-Sen's ARAA had this already).
    \item Oversmoothing and how to avoid it (see recent Bronstein papers)
    \item Graph Neural Networks as the natural interface for Symbolic Regression
    \item "Physics is a inherently a relational science"
    \item The text should be extremely analogy-heavy
\end{itemize}


There is no page limit - but aiming for 15-20 would be good